\chapter{Conclusion and Outlook}
\label{ch:conclusion}

\section{Synthesis of the research findings}
\label{sec:assessment_observed_findings}

The central research problem was how to use rapid activated sludge predictions for engineering analysis when prediction accuracy alone does not establish a usable treatment state or a good operating decision. The dissertation addressed this problem through complete component prediction, physical enforcement, explicit regression, and independent assessment of selected controls. Its general answer is that these capabilities can be connected through a common component-state representation. They still require separate evidence. A physically consistent state can be inaccurate, an inspectable regression can have higher error than a flexible learner, and a faster search can select a worse decision on the original mechanistic basis.

The adopted component state retained organic substrate, nitrogen species, phosphorus storage, biomass, and solids. The reporting map then produced COD, TN, TP, and TSS \citep{Henze2006}. This order allowed the same response to support material accounting, treatment reporting, and operating assessment without treating a reported total as a unique component state. The findings below answer the four research questions in Section~\ref{sec:research_questions}. Their numerical evidence is reported in Sections~\ref{sec:projection_results}, \ref{sec:icsor_results}, and \ref{sec:plant_results}.

\subsection{Complete reactor prediction and data availability}

The first research question concerned the accuracy of different learners as training data changed and inputs moved beyond the sampled domain. The common reactor benchmark showed that the preferred learner depended on the simulation budget. At 500 cases, XGBoost and LightGBM had the lowest observed raw component nRMSE values of 0.435 and 0.439. At 10,000 cases, the multilayer perceptron had the lowest raw and projected values of 0.150 and 0.149. This answers the comparative gap with evidence from a shared component target rather than separate prediction tasks. It also shows why selection at one training size cannot establish the preferred model at every size \citep{VieringLoog2023,Borisov2024}.

The exterior assessment supplied a different answer about the usable input range. The multilayer perceptron again had the lowest error, but its projected nRMSE was 0.251 under mild exterior inputs and 0.521 under severe exterior inputs. Every projected exterior state still passed the physical checks. The larger severe-domain error therefore did not arise from failure of those checks. For practical use, an unfamiliar loading or operating condition needs an accuracy assessment even when projection returns a valid balance. The benchmark supports model selection within declared data and domain conditions, not unrestricted extrapolation or one universal learner ranking.

\subsection{Physical enforcement and its consequences}

The second research question concerned what projection changes in the same raw prediction. Every raw assessment state failed the specified mass conservation tolerance. Some predictors returned only non-negative raw concentrations, which showed that positivity alone was insufficient. Joint enforcement returned states that passed both requirements in all 766,350 assessed projection instances. These instances included repeated assessment across nested design sizes rather than that many distinct simulated reactor states. The maximum invariant residuals were $2.06\times10^{-13}$ in the learning domain and $1.72\times10^{-13}$ in the exterior assessment. The result demonstrates numerical enforcement of the adopted balances and bounds across different learners. It applies the existing projection of \citet{KircherVotsmeier2025} to a common activated sludge task rather than introducing a new projection principle.

The accuracy effect was mixed. At the largest design size, projection improved 188 of 260 predictor--component comparisons, but aggregate component nRMSE decreased for only five of the thirteen predictors. Every predictor had lower TN and TP root error and higher COD and TSS root error. Thus, physical enforcement closed the compliance gap without eliminating the approximation gap. A treatment analysis concerned mainly with nutrients can reach a different assessment of the adjustment from one concerned with organic matter, biomass, or solids. The result supports separate component and composite reporting and confirms the importance of the adjustment metric discussed by \citet{SturmSilva2024}.

Displacement and computational cost completed this answer. They described how much a raw prediction changed and the effort needed to return the checked state. Projection dominated the repeated prediction cost of the fastest raw linear and neural models under the stated batch protocol. A rapid regression evaluation is therefore not the same as a rapid complete prediction. The paired assessment provides a basis for selecting an enforcement procedure together with a learner while retaining accuracy, physical compliance, and complete prediction cost as distinct criteria \citep{BhosekarIerapetritou2018}.

\subsection{Interpretation of the complete regression response}

The third research question concerned how an explicit regression could support interpretable prediction with mass conservation and non-negativity. ICSOR supplied named operating, loading, curvature, and interaction terms, followed by learned component coupling and checked projection \citep{Selerio2026}. These stages made it possible to trace an input contribution to the returned component state. In one fitted model, 44 of 210 unique COD influent-curvature terms exceeded their reporting threshold, while 371 of 380 shared off-diagonal coupling entries exceeded theirs. The representation was selective in influent curvature but dense in coupling. Its contribution to interpretation was an explicit response structure, not a uniformly sparse model or uniquely identified biological parameters.

The accuracy evidence supported a conditional practical use for this structure. At 500 cases, ICSOR had the lowest mean held-out pooled composite root error of 10.79, compared with 11.26 for XGBoost and 11.45 for LightGBM. It also had the lowest reported root-error learning-curve area of 6.33 after normalization by the design-size interval. At 10,000 cases, repeated-split errors were 5.94 for ICSOR and 4.62 for the multilayer perceptron. LightGBM had the best mean rank across the broader assessment. These outcomes favor considering ICSOR when an explicit equation and limited simulation data matter, while retaining flexible learners when the main priority is lower full-size error. Accuracy scoring used checked ICSOR states and affine-aligned comparator states, so these comparisons do not isolate the effect of regression structure from output enforcement.

The physical assessment showed that interpretation must include projection. Affine projection sufficed for 436 of the 2000 fixed-partition cases. The remaining 1564 cases required constrained linear projection, and every final state passed both adopted physical checks. The analytical raw and affine sensitivities and the hypothetical example of an active concentration bound showed why a raw coefficient is not necessarily the final local slope. An engineer examining an aeration response must follow coupling and the active projection constraints at the operating point of interest. The model therefore addresses the interpretation gap through a traceable calculation, while coefficient stability and causal process identification remain separate questions \citep{Selerio2026,Shen2023}.

\subsection{Connected prediction and operating decisions}

The fourth research question concerned the support that joint prediction and projection could provide for connected plant decisions. The response retained the mixer, five biological reactors, clarifier overflow, clarifier underflow, and stored solids. Joint enforcement linked shared component flows and inventories rather than checking each unit in isolation \citep{Alex2008,JeppssonDiehl1996}. The development and holdout designs retained 7386 and 1845 mechanistic states. On the holdout, projection reduced aggregate nRMSE from 0.04690 to 0.04306, an improvement of 8.19\%, and reduced aggregate nMAE by 15.65\%. It improved 30 of 32 location--composite nRMSE values. The two increases were COD in the final reactor and overflow. This evidence addresses the network-consistency gap and shows a broad, but not uniform, statistical benefit from the complete connected procedure.

The decision assessment identified a limit that the holdout improvement could not resolve. In S4, Extended ICSOR predicted TN and TP concentrations of 2.67 and 2.08 mg L$^{-1}$, while independent mechanistic evaluation gave 7.47 and 4.68 mg L$^{-1}$. These errors remained after projection and trust screening. A nutrient requirement between the predicted and reference concentrations would lead to different judgments about the selected setting. The result answers the decision-quality gap by showing why treatment outcomes must be calculated at the controls actually selected, not inferred from average holdout performance \citep{Alexandrov1998,Pedrozo2025}.

Smooth NLP had a lower common-reference objective in the nominal case and all ten influent scenarios. Nominal values were 0.8647 for Extended ICSOR and 0.7862 for Smooth NLP. Extended ICSOR had shorter recorded search times but did not match the direct route's objective values. The answer to the fourth question is therefore qualified. Joint prediction and projection can support rapid exploration with a consistent material account, but they do not replace mechanistic decision evaluation. The comparisons concern available local decisions under the declared priorities and do not establish a global optimum or superiority of one route in every plant setting.

\section{Contributions and gaps addressed}
\label{sec:assessment_conclusions}

The dissertation contributes a common complete-component reactor benchmark. This addresses the lack of directly comparable evidence about learner behavior under changing data availability and input conditions. It retains both small component concentrations and reported treatment quantities. The contribution is the defined activated sludge comparison and its evidence, not a claim that one learner is always preferable.

A second contribution is the paired assessment of a common physical projection across learners. It makes the distinction between physical enforcement and statistical improvement measurable. The adopted projection and stoichiometric foundations come from existing work \citep{SturmWexler2022,KircherVotsmeier2025}. The dissertation's contribution lies in applying them to complete activated sludge predictions and evaluating compliance, error redistribution, displacement, and cost together. This also establishes the limits of interpreting a feasible component state as a fully accurate mechanistic response.

A third contribution is the structured regression and analysis of its complete returned response. Named terms expose operating and loading relationships, while component coupling represents shared output dependence. Stage-specific physical diagnostics and sensitivities distinguish what the regression learns from what projection imposes. This addresses the gap between having visible coefficients and explaining the final treatment prediction. The structure supports inspection without presenting fitted associations as causal kinetic parameters \citep{Selerio2026}.

The fourth contribution is a connected component-state formulation with joint projection and independent decision verification. It includes the shared streams and solids inventories that a reactor-only output cannot represent. Common mechanistic evaluation then connects those predictions to treatment and resource choices. The contribution includes the observed limitation of the surrogate decisions. Demonstrating lower search time alongside higher reference objectives supplies useful decision evidence rather than assuming that computational speed establishes engineering superiority.

Table~\ref{tab:integrated_observed_findings} links the gaps to their answers and practical uses. The integration establishes a consistent assessment logic, not one numerical ranking across all three studies. Reactor component errors use training-based scales, the structured comparison pools native composite coordinates, and plant holdout errors use location-specific reference ranges. Their validation designs and enforcement procedures also differ. These scores should therefore be interpreted within their own assessments rather than combined into one cross-study measure.

\begin{table}[htbp]
\centering\small
\caption{Research gaps, observed answers, and practical uses of the integrated findings.}
\label{tab:integrated_observed_findings}
\begin{tabular}{>{\raggedright\arraybackslash}p{0.24\textwidth}>{\raggedright\arraybackslash}p{0.40\textwidth}>{\raggedright\arraybackslash}p{0.26\textwidth}}
\toprule
Research gap & Observed answer & Practical use\\
\midrule
Uncertain complete-state learner performance & The leading learner changed with data budget, and exterior accuracy deteriorated despite physical compliance & Select for the available data and intended input range\\
Training does not establish physical compliance & All 766,350 projected instances passed both checks, but aggregate error improved for only five of thirteen learners at the largest size & Check balances and bounds separately from treatment accuracy\\
Visible terms do not explain the final response & ICSOR exposed operating and loading terms, while most fixed-partition cases required constrained projection & Interpret the coupled and projected response at the operating point\\
Consistent unit predictions do not establish good plant decisions & Joint prediction improved 30 of 32 local scores, but faster surrogate searches selected higher reference objectives & Verify selected controls with the original plant equations\\
\bottomrule
\end{tabular}
\end{table}

Together, these contributions connect process knowledge with statistical approximation and decision evidence. Reaction stoichiometry supplies the invariant relations, while unit boundaries supply transport and inventory relations \citep{Henze2006,Alex2008}. Regression approximates the response within the sampled conditions. Projection imposes the declared physical requirements. Interpretation examines the complete mapping, and independent mechanistic evaluation establishes the reference outcome at selected controls. The integrated contribution is this defined connection and the evidence needed to assess each part. It does not claim invention of activated sludge stoichiometry, general projection, or surrogate optimization.

\section{Practical implications}
\label{sec:assessment_practical_implications}

\subsection{Treatment and resource priorities}

The operating results help engineers compare treatment gains with the resources required to obtain them. In S10, the direct decision reduced normalized quality cost from 1.4748 to 0.3699 while increasing the weighted resource contribution from 0.0557 to 0.4191. The total objective decreased from 0.7931 to 0.6040. This choice accepted greater capacity and resource contributions for better weighted treatment. In S8, a lower resource contribution produced a slightly lower total objective while all four effluent composites were slightly higher. Reporting the separate contributions helps engineers understand the selected setting and relate it to their treatment priorities \citep{PadronPaez2020,Aboagye2021}.

The selected HRT values also support capacity planning. Varying HRT at fixed fresh flow changed the biological reactor volume. The direct route reached the 36-h upper bound in S5 and S10, while the surrogate reached the 6-h lower bound in several cases. Comparing these settings connects reactor capacity with the treatment obtained. Aeration can be assessed alongside reactor volume because both influence oxygen-transfer capacity. Wasted-solids production can likewise be assessed from wasting flow and underflow concentration. For operating analysis at an installed plant, the existing reactor volumes and equipment capacities provide the basis for comparing control settings \citep{Henze2008,Miederer2025}.

The separation of quality and resource terms provides a practical basis for developing facility-specific objectives. Effluent concentrations and removal values can express treatment targets. Electricity demand, recycle pumping, and sludge-handling costs can express resource priorities. Their weights specify how each quantity contributes to the preferred decision. Engineers can use the resulting contributions to explain a choice and compare alternatives under different priorities \citep{PadronPaez2020,Miederer2025}.

\subsection{A verified decision workflow}

The surrogate supports screening and candidate selection. The engineer first defines the plant boundary, influent composition, available controls, and treatment targets. The fitted model then provides component responses for comparing alternatives. Joint projection reconciles the predicted streams and inventories before treatment and resource quantities are compared. Physical residuals and projection displacement provide diagnostic information for selecting candidates for detailed evaluation.

Independent evaluation with the original reaction and settling equations supplies final-effluent concentrations, shared solids flows, inventories, and the common-reference objective at the selected controls. These quantities provide a consistent basis for choosing among candidate settings. The evaluations can also guide additional mechanistic cases near promising operating conditions. This connects further surrogate fitting with the decisions that the model is intended to support \citep{Alexandrov1998,Hameed2026}.

Component and location information also supports process diagnosis. The increase from fresh-influent to mixer concentrations shows the material circulated through recycle. The TSS decrease across the clarifier shows the contribution of solids separation. Separate ammonium and oxidized-nitrogen concentrations reveal changes among nitrogen forms when TN changes little. Examining these quantities alongside returned solids and inventories helps engineers relate aeration, recycle, and wasting choices to the treatment response \citep{Henze2006,Alex2008,JeppssonDiehl1996}.

\subsection{Computational value}

The measured scenario mean initial-search times were 1.726 s for Extended ICSOR and 5.236 s for Smooth NLP. Their ratio was approximately 3.03. The shorter surrogate search interval provides computational value for repeated exploration of operating alternatives. Rapid screening is useful when an engineer needs to compare several influent conditions or combinations of aeration, recycle, and wasting settings.

Repeated application distributes model-development effort across multiple operating comparisons. Engineers can use a fitted surrogate to explore alternatives and reserve detailed mechanistic calculations for selected candidates. This workflow focuses process calculations on the operating decisions of interest while using rapid prediction for broader comparison \citep{BhosekarIerapetritou2018,Bliek2023}.

\section{Directions for future research}

\subsection{Adaptive sampling around operating decisions}

Future studies can generate new mechanistic cases where the surrogate is most influential in the operating search. A first surrogate can screen the declared domain and identify candidate controls. Mechanistic calculations at those candidates can then supply additional training states. Sampling can concentrate on large nutrient discrepancies, strong projection displacements, and regions close to treatment requirements. Repeating fitting and operating assessment would connect the simulation budget to the decisions it is intended to support. This proposal develops the reference-informed model management of \citet{Alexandrov1998} and the gray-box approximation management discussed by \citet{Hameed2026}.

An informative comparison would hold the mechanistic evaluation budget fixed and compare uniform sampling, error-based sampling, and decision-based sampling. The outcomes should include effluent errors at selected controls, the common-reference objective, the number of accepted decisions, and the complete time required to deliver them. The S4 nutrient discrepancies provide a concrete target for such a study. Additional cases near those selected operating conditions could test whether focused sampling improves TN and TP prediction where the operating choice depends on them.

\subsection{Projection aligned with treatment priorities}

The reactor results motivate further study of projection weights and the quantities they emphasize. A projection that minimizes changes in component coordinates can redistribute error differently from an assessment based on reported nutrients or solids. Future formulations can compare component scaling, composite-sensitive weighting, and weights based on prediction uncertainty. Each formulation should retain the same mass conservation and non-negativity requirements so that changes in treatment accuracy can be attributed to the adjustment metric. The species-weighted correction of \citet{SturmSilva2024} provides a precedent for choosing adjustment weights from component magnitude and uncertainty.

The connected plant offers an additional test of this idea. Weighting can distinguish internal state accuracy from final-effluent accuracy while preserving the unit connections. Studies can assess whether nutrient-focused weights improve selected-control predictions without producing large errors in biomass, solids inventory, or other treatment locations. This would make the relation between physical adjustment and the intended engineering decision more explicit.

\subsection{Uncertainty-aware prediction and operating safeguards}

Uncertainty can be developed alongside the projected point prediction. Repeated fits, model ensembles, or probabilistic component predictions can describe variation in the estimated treatment response. Each sampled state can be projected onto the same physical set. The resulting distributions can then be evaluated for their coverage of mechanistic reference concentrations, particularly for TN, TP, and overflow TSS. The conservation-aware probabilistic framework of \citet{Hansen2024} and the differentiable probabilistic projection of \citet{Utkarsh2025} provide starting points for this extension.

Operating studies can use these distributions to formulate treatment safeguards. For example, a candidate setting can be selected to keep an upper predictive bound for TP below a specified treatment requirement. Mechanistic and plant-based evaluations can test the frequency with which those safeguards are met. The effluent ammonium safeguards considered by \citet{Ren2026} illustrate how treatment requirements can guide aeration decisions. Separating uncertainty associated with surrogate fitting, influent measurements, and mechanistic parameters would also show which information most improves the decision \citep{Machado2014,Petersen2003}.

\subsection{Stable interpretation of the complete response}

Further interpretation work can examine the stability of coefficient groups across repeated training samples and operating domains. The present structure provides natural groups for operating terms, loading terms, curvature, and coupling. Resampling can show which groups remain prominent and which individual coefficients vary. Comparisons based on the complete predicted response can help identify different coefficient combinations that describe similar treatment behavior. This work would address the coefficient-stability question identified by \citet{Selerio2026}.

Sensitivity studies should also follow the final projection. Derivatives can be evaluated at representative operating points and checked against direct perturbations of the returned concentrations. Mapping the active constraints would show where a nutrient response changes smoothly and where a projection boundary alters its slope. Comparing these projected sensitivities with mechanistic sensitivities would connect model inspection more closely to aeration, recycle, and wasting decisions \citep{Selerio2026,Shen2023}.

\subsection{Dynamic prediction and control}

A dynamic extension can replace steady-state targets with trajectories of component concentrations and stored solids. The governing balances would include the time derivatives of reactor and clarifier inventories \citep{Takacs1991,Alex2008}. Training data can represent changing influent loads, aeration schedules, recycle settings, and sludge wasting. Physical enforcement can then account for the material transported and transformed during each prediction interval.

This extension would support studies of predictive control under daily influent variation and load disturbances. The evaluation can include transient nutrient peaks, recovery time, inventory changes, oxygen demand, and the computation required at each control update. The field-tested aeration controller of \citet{Bernardelli2020} and the wastewater control review of \citet{Croll2023} provide examples for developing and evaluating such dynamic strategies. The connected component representation provides a useful starting point because it already retains the locations and inventories that influence delayed treatment responses.

\subsection{Plant validation, topology, and resource models}

Validation with measured plant data is a major next step. Studies can combine influent fractionation, calibrated mechanistic parameters, unit concentrations, clarifier solids measurements, and recorded controls \citep{Petersen2003,Machado2014,Newhart2019}. A prospective assessment can fix the fitting and operating procedure before collecting the final evaluation period \citep{Kaufman2012}. Testing across seasons and loading conditions would provide evidence about how the component surrogate transfers to observed treatment behavior.

Operating objectives can also be developed for specific facilities. An installed plant has defined reactor volumes, aeration capacity, recycle capacity, and discharge requirements. Electricity demand, chemical use, sludge handling, and operating costs can replace or supplement the engineering proxies used here \citep{Henze2008,Miederer2025}. Multiobjective studies can present the attainable treatment and resource combinations before a preference is selected. Extensions to other reactor arrangements, attached-growth units, and resource-recovery connections can retain the same principle of defining component states and boundary streams before physical enforcement is imposed \citep{Henze2008,Aboagye2021,Durkin2024}.

More extensive search comparisons can evaluate several initial controls, hybrid surrogate and mechanistic searches, and the fifteen-start basin-sensitivity design described in Chapter~\ref{ch:plant}. Recording state generation, fitting, prediction, projection, search, and verification costs would establish when surrogate development becomes worthwhile for repeated decisions \citep{BhosekarIerapetritou2018,Bliek2023}. The central assessment would remain the treatment outcome and common-reference objective obtained for a given computational budget.

\section{Concluding perspective}

The dissertation shows how complete component prediction, physical enforcement, interpretation, and decision verification can form one activated sludge analysis framework. Its findings answer the research questions without treating one capability as proof of another. Learner choice depended on the available data and input range. Projection enforced the adopted physical requirements but did not improve every error. Structured regression made the response inspectable but did not have the lowest error at the largest data budget. Connected projection improved plant prediction, while mechanistic evaluation exposed higher objectives at the faster surrogate-selected decisions.

The practical contribution is a defined basis for rapid exploration with explicit checks on what is predicted and what is selected. The surrogate can help screen alternatives, while the original mechanistic model supplies the reference treatment outcome at candidate controls. Plant measurements and facility-specific priorities are needed to extend that basis to actual operation. Computational efficiency is useful when it helps deliver an acceptable treatment decision, not when search time alone is reduced.
